pdfoutput=1
\pdfoutput=1
%=============================================================================
% Supplementary Information S2
% Theorem 2: Weak Feature Failure Boundary
% Nature SI — LaTeX
%=============================================================================

\section{Theorem 2: Weak Feature Failure Boundary (The Impossible Demand Theorem)}
\label{sec:thm2}

When the feature representation $\phi(x)$ contains insufficient information
about the true state $S$, consistency-based noise detection methods
(including SCX) cannot outperform the loss-based baseline. Theorem~2
quantifies this boundary from an information-theoretic perspective.

%------------------------------------------------------------------------------
\subsection{Information-Theoretic Setup}
\label{sec:thm2:setup}

\paragraph{Data generation.}
Let $X \in \mathcal{X} \subseteq \mathbb{R}^d$ be the input, $Y \in \mathcal{Y}$
the label, and $Z \in \{0,1\}$ the noise indicator ($Z=1$ for noisy samples).
The true (unobserved) state is $S = s(X) \in \mathcal{S}$ with
$|\mathcal{S}| = K_{\mathcal{S}}$. The observed feature representation is
$\phi(X) \in \Phi \subseteq \mathbb{R}^{d_\phi}$.

SCX assumes the following generative structure:
\[
Z \;\to\; S \;\to\; (X, Y) \;\to\; \phi(X),
\qquad
Z \;\to\; S \;\to\; X \;\to\; \{f_m(X)\}.
\]
This yields the Markov chain
\begin{equation}
Z \to S \to X \to \phi(X),
\label{eq:markov-chain}
\end{equation}
and by the data processing inequality,
\[
I(\phi(X); Z) \leq I(X; Z) \leq I(S; Z) \leq H(Z) \leq \log 2.
\]

\paragraph{Weak feature definition.}

\begin{definition}[$\delta$-weak feature]
\label{def:weak-feature}
A feature map $\phi: \mathcal{X} \to \Phi$ is called $\delta$-weak with respect
to the true state map $s: \mathcal{X} \to \mathcal{S}$ if
\begin{equation}
I(\phi(X); s(X)) \leq \delta,
\label{eq:weak-feature}
\end{equation}
where $I(\cdot;\cdot)$ is the Shannon mutual information, measured in nats.
\end{definition}

The normalised weakness is
\begin{equation}
\varepsilon_\phi = \frac{\delta}{\log K_{\mathcal{S}}} \in [0, 1],
\label{eq:norm-weakness}
\end{equation}
since $I(\phi; S) \leq H(S) \leq \log K_{\mathcal{S}}$.
$\varepsilon_\phi = 1$ corresponds to features with no state information
(worst case), $\varepsilon_\phi = 0$ to full state information.

An equivalent definition uses the total variation distance. For two states
$s_1, s_2$, let
$\operatorname{TV}_{s_1,s_2} = \operatorname{TV}(P_{\phi \mid S=s_1},
P_{\phi \mid S=s_2})$. Then the feature weakness is
\begin{equation}
\Delta_\phi = \max_{s_1 \neq s_2}
\operatorname{TV}(P_{\phi \mid S=s_1}, P_{\phi \mid S=s_2}),
\label{eq:tv-weakness}
\end{equation}
and by Pinsker's inequality,
\begin{equation}
\Delta_\phi \leq \sqrt{\frac{\delta}{2}}.
\label{eq:pinsker-tv}
\end{equation}
\textbf{[Heuristic bound.]} This inequality is not rigorous as stated:
Pinsker's inequality bounds $\operatorname{TV}(P,Q) \leq \sqrt{\operatorname{KL}(P\|Q)/2}$,
but $\delta = I(\phi;S)$ is the mutual information, which equals
$\operatorname{KL}(P_{\phi,S}\|P_\phi P_S)$, not a KL divergence between two
state-conditional distributions. The bound $\Delta_\phi \leq \sqrt{\delta/2}$
holds approximately when state probabilities are balanced and the
state-conditional distributions are well-separated, but it should be treated
as a heuristic rather than a formal consequence of Pinsker. A rigorous
bound would require additional factors depending on $\min_s \mathbb{P}(S=s)$.
When $\Delta_\phi$ is small, points from different true states are
indistinguishable in $\phi$-space.

\paragraph{Loss baseline detector.}
The loss-based detector uses no state information:
\[
\hat{z}_{\text{loss}}(x) = \mathbf{1}\{\max_m \ell(f_m(x), y) > t\},
\]
where $\{f_m\}_{m=1}^M$ are $M$ expert models. Denote its optimal performance
by $\operatorname{AUC}_{\text{base}}$,
$\operatorname{PRAUC}_{\text{base}}$, and $\operatorname{F1}_{\text{base}}$.

For a completely random detector, the optimal performance is
\begin{equation}
\operatorname{AUC}_{\text{rand}} = 0.5,\qquad
\operatorname{PRAUC}_{\text{rand}} = \eta,\qquad
\operatorname{F1}_{\text{rand}} = \frac{2\eta}{1+\eta},
\label{eq:random-baseline}
\end{equation}
where $\eta = \mathbb{P}(Z=1)$ is the marginal noise rate.

%------------------------------------------------------------------------------
\subsection{Lemma 1: State Estimation Error (Fano)}
\label{sec:thm2:lemma1}

\begin{lemma}[State estimation error lower bound -- Fano's inequality]
\label{lem:fano}
Let $\hat{S}$ be any estimator of the true state $S$ based on $\phi(X)$
(i.e., the output of the SCX state-discovery algorithm). If $\phi$ is
$\delta$-weak, then
\begin{equation}
\mathbb{P}(\hat{S} \neq S) \;\geq\;
\frac{H(S) - \delta - \log 2}{\log K_{\mathcal{S}}}.
\label{eq:fano-bound}
\end{equation}
\end{lemma}

\begin{proof}
Apply Fano's inequality \cite{fano1961,cover2006}:
\begin{align*}
\mathbb{P}(\hat{S} \neq S)
&\geq \frac{H(S \mid \phi(X)) - \log 2}{\log K_{\mathcal{S}}} \\
&= \frac{H(S) - I(\phi(X); S) - \log 2}{\log K_{\mathcal{S}}} \\
&= \frac{H(S) - \delta - \log 2}{\log K_{\mathcal{S}}}.
\qedhere
\end{align*}
\end{proof}

\begin{corollary}[Uniform states]
\label{cor:fano-uniform}
If the true states are approximately uniform,
$H(S) \approx \log K_{\mathcal{S}}$, then
\[
\mathbb{P}(\hat{S} \neq S) \;\geq\; 1 - \frac{\delta + \log 2}{\log K_{\mathcal{S}}}.
\]
As $\delta \to 0$ with $K_{\mathcal{S}}$ fixed,
$\mathbb{P}(\hat{S} \neq S) \to 1 - \frac{\log 2}{\log K_{\mathcal{S}}} > 0$.
For large $K_{\mathcal{S}}$ this limit approaches $1$.
\end{corollary}

\begin{corollary}[Two-state case]
\label{cor:fano-two-state}
When $K_{\mathcal{S}} = 2$,
\[
\mathbb{P}(\hat{S} \neq S) \;\geq\;
\frac{H(S) - \delta - \log 2}{\log 2}.
\]
\end{corollary}

\begin{remark}[Achievability]
\label{remark:fano-achievability}
The converse of Fano's inequality guarantees the \emph{existence} of an
estimator $\hat{S}$ operating on $\phi$ such that
\[
\mathbb{P}(\hat{S} \neq S) \;\leq\; \frac{\delta + \log 2}{\log K_{\mathcal{S}}}.
\]
This existence bound does not guarantee that a particular algorithm
(e.g.\ $k$-means) attains it. In practice, the SCX clustering error satisfies
\[
\mathbb{P}(\hat{S} \neq S) \;\leq\; O\!\left(\frac{\delta}{\log K_{\mathcal{S}}}\right)
+ \varepsilon_{k\text{-means}},
\]
where $\varepsilon_{k\text{-means}}$ captures the approximation error of
$k$-means (initialisation, non-convexity, non-spherical clusters). When
$\delta$ is small (weak features), \emph{any} algorithm is fundamentally
limited by the Fano lower bound.
\end{remark}

%------------------------------------------------------------------------------
\subsection{Lemma 2: SCX Reliability Degradation}
\label{sec:thm2:lemma2}

\begin{lemma}[SCX reliability estimation degradation]
\label{lem:scx-degradation}
Let $C(s)$ be the consensus score for true state $s$, and let $C(\hat{s})$
be the consensus score for an estimated state $\hat{s}$ obtained by clustering
$\phi$. If $\phi$ is $\delta$-weak, then
\begin{equation}
\mathbb{E}\bigl[|C(\hat{S}) - C(S)|\bigr]
\;\leq\; \mathbb{P}(\hat{S} \neq S) + O\!\left(\frac{1}{\sqrt{n_{\min}}}\right),
\label{eq:degradation-bound}
\end{equation}
where $n_{\min} = \min_{s \in \mathcal{S}} n_s$ is the minimum per-state
sample size.
\end{lemma}

\begin{proof}
Decompose the expectation into correctly and incorrectly estimated states:
\begin{align*}
\mathbb{E}\bigl[|C(\hat{S}) - C(S)|\bigr]
&= \mathbb{P}(\hat{S} = S)\,
   \mathbb{E}\bigl[|C(\hat{S}) - C(S)| \mid \hat{S} = S\bigr] \\
&\qquad + \mathbb{P}(\hat{S} \neq S)\,
   \mathbb{E}\bigl[|C(\hat{S}) - C(S)| \mid \hat{S} \neq S\bigr].
\end{align*}

\emph{Case 1 ($\hat{S} = S$):} When the state is correctly estimated,
$C(\hat{S})$ is a consistent estimator of $C(S)$. By Chebyshev's inequality
and $C(\cdot) \in [0,1]$,
\[
\mathbb{E}\bigl[|C(\hat{S}) - C(S)| \mid \hat{S} = S\bigr]
\;\leq\; \frac{\sigma_C}{\sqrt{n_s}}
\;\leq\; O\!\left(\frac{1}{\sqrt{n_{\min}}}\right),
\]
where $\sigma_C$ is the standard deviation of $C$.

\emph{Case 2 ($\hat{S} \neq S$):} When the state estimate is incorrect,
$C(\hat{S})$ is a mixture of the consensus scores of the true states:
\[
C(\hat{S}) = \sum_{s \in \mathcal{S}} \mathbb{P}(S = s \mid \hat{S})\, C(s).
\]
Since $C(\cdot) \in [0,1]$, the pointwise difference satisfies
$|C(\hat{S}) - C(S)| \leq 1$, and consequently
$\mathbb{E}[|C(\hat{S}) - C(S)| \mid \hat{S} \neq S] \leq 1$.

Combining the two cases,
\begin{align*}
\mathbb{E}\bigl[|C(\hat{S}) - C(S)|\bigr]
&\leq \mathbb{P}(\hat{S}=S) \cdot O(1/\sqrt{n_{\min}})
   + \mathbb{P}(\hat{S}\neq S) \cdot 1 \\
&\leq \mathbb{P}(\hat{S} \neq S) + O(1/\sqrt{n_{\min}}). \qedhere
\end{align*}
\end{proof}

\begin{corollary}[Degradation to global consistency]
\label{cor:global-consistency}
As $\delta \to 0$, Lemma~\ref{lem:fano} gives
$\mathbb{P}(\hat{S} \neq S) \to 1 - \frac{\log 2}{\log K_{\mathcal{S}}}$.
In this limit,
\[
C(\hat{S}) \xrightarrow{p} \sum_{s \in \mathcal{S}} \rho(s)\, C(s) = \bar{C},
\]
where $\rho(s) = \mathbb{P}(S = s)$ and $\bar{C}$ is the global average
consistency. When features carry no state information, all estimated states
converge to the same consistency value $\bar{C}$.
\end{corollary}

\begin{corollary}[Noise score degradation]
\label{cor:ns-degradation}
Under the same conditions and the additional assumption that the estimated
states are approximately balanced
($\max_{\hat{s}} \rho(\hat{s}) / \min_{\hat{s}} \rho(\hat{s}) \leq R$ for
some constant $R$), the SCX noise score degenerates to
\[
\operatorname{NS}(x) \propto r(x),
\]
where $r(x)$ is the residual of sample $x$. The SCX detector therefore
reduces to the loss-based baseline detector. If the estimated states are
highly imbalanced, the degradation is approximate with an additional
$O(1/\sqrt{n})$ bias from state proportion estimation error.
\end{corollary}

%------------------------------------------------------------------------------
\subsection{Proof of Theorem 2: AUC, PR-AUC, and F1 Bounds}
\label{sec:thm2:proof}

\begin{theorem}[Weak feature failure lower bound]
\label{thm:weak-feature}
Let $\phi: \mathcal{X} \to \Phi$ be a $\delta$-weak feature map with respect to
the true state $S$, i.e.\ $I(\phi(X); S) \leq \delta$ (in nats). Let
$h_{\text{SCX}}$ be the SCX noise detector operating under the standard
pipeline (clustering on $\phi$, state-conditional reliability estimates,
noise scoring). Assume the estimated states are approximately balanced
($\max_{\hat{s}} \rho(\hat{s}) / \min_{\hat{s}} \rho(\hat{s}) \leq R$) and
the clustering algorithm does not introduce error beyond the Fano lower bound.
Then:

\begin{enumerate}
\item[(a)] \textbf{AUC bound:}
\begin{equation}
\boxed{\;
\operatorname{AUC}(h_{\text{SCX}}) \;\leq\;
\operatorname{AUC}_{\text{base}} +
\sqrt{\frac{\delta}{2}}\cdot\left(\frac{1}{\eta} + \frac{1}{1-\eta}\right)
\;}
\label{eq:auc-bound}
\end{equation}

\item[(b)] \textbf{PR-AUC bound:}
\begin{equation}
\boxed{\;
\operatorname{PRAUC}(h_{\text{SCX}}) \;\leq\;
\operatorname{PRAUC}_{\text{base}} +
\sqrt{\frac{\delta}{2}}\cdot\left(\frac{1}{\eta} + \frac{1}{1-\eta}\right)
\;}
\label{eq:prauc-bound}
\end{equation}

\item[(c)] \textbf{F1 bound:}
Assume the detector operates in a regime where both precision and recall
are bounded below by $\varepsilon > 0$ (typically $\varepsilon = 0.1$).
Then there exists a constant $C_F = 2/\varepsilon^2$ such that
\begin{equation}
\boxed{\;
\operatorname{F1}(h_{\text{SCX}}) \;\leq\;
\operatorname{F1}_{\text{base}} + C_F \cdot \sqrt{\frac{\delta}{2}}
\;}
\label{eq:f1-bound}
\end{equation}
\end{enumerate}

Here $\operatorname{AUC}_{\text{base}}$,
$\operatorname{PRAUC}_{\text{base}}$, and $\operatorname{F1}_{\text{base}}$
are the performance metrics of the loss-based baseline detector that thresholds
$\max_m \ell(f_m(x), y)$ without using $\phi$ or state information.
\end{theorem}

\begin{proof}
The proof proceeds in six steps.

\paragraph{Step 1: Construct an auxiliary distribution $\tilde{P}$.}
Define $\tilde{P}$ by forcing $\phi$ and $S$ to be independent while preserving
their marginal distributions:
\begin{equation}
\tilde{P}(\phi, S) = P(\phi)\, P(S).
\label{eq:aux-dist}
\end{equation}
All other conditional distributions (given $S$, the distribution of $(X,Y)$,
experts, etc.) are unchanged.

\paragraph{Step 2: Relate KL divergence and mutual information.}
The KL divergence between $P$ and $\tilde{P}$ equals the mutual information
between $\phi$ and $S$:
\begin{align*}
\operatorname{KL}(P \parallel \tilde{P})
&= \iint P(\phi, S) \log\frac{P(\phi, S)}{P(\phi)P(S)}\, d\phi\, dS
= I(\phi(X); S) = \delta.
\end{align*}
By Pinsker's inequality \cite{pinsker1964},
\begin{equation}
\operatorname{TV}(P, \tilde{P}) \;\leq\;
\sqrt{\frac{\operatorname{KL}(P \parallel \tilde{P})}{2}}
= \sqrt{\frac{\delta}{2}}.
\label{eq:pinsker-tv-bound}
\end{equation}

\paragraph{Step 3: Transfer the TV bound to the prediction distribution.}
Let $P_{\text{pred}}$ and $\tilde{P}_{\text{pred}}$ be the joint distributions
of $(\hat{z}_{\text{SCX}}(X), Z)$ under $P$ and $\tilde{P}$, respectively.
Since the SCX noise score is a deterministic function of $X$ (and hence of
$\phi$ and $S$), the data processing inequality \cite{cover2006} yields
\begin{equation}
\operatorname{TV}(P_{\text{pred}}, \tilde{P}_{\text{pred}})
\;\leq\; \operatorname{TV}(P, \tilde{P})
\;\leq\; \sqrt{\frac{\delta}{2}}.
\label{eq:dp-tv}
\end{equation}

\paragraph{Step 4: Analyse SCX under $\tilde{P}$.}
Under $\tilde{P}$, $\phi \perp S$, so $\phi$ carries no information about $S$.
By Corollary~\ref{cor:ns-degradation}, the SCX noise score degenerates to the
residual:
\[
\operatorname{NS}_{\tilde{P}}(x) \propto \max_m \ell(f_m(x), y).
\]
Hence the SCX detector under $\tilde{P}$ is equivalent to the loss-based
baseline detector:
\begin{align}
\operatorname{AUC}_{\tilde{P}}(h_{\text{SCX}}) &= \operatorname{AUC}_{\text{base}},
\label{eq:base-equiv-au} \\
\operatorname{PRAUC}_{\tilde{P}}(h_{\text{SCX}}) &= \operatorname{PRAUC}_{\text{base}},
\label{eq:base-equiv-prau} \\
\operatorname{F1}_{\tilde{P}}(h_{\text{SCX}}) &= \operatorname{F1}_{\text{base}}.
\label{eq:base-equiv-f1}
\end{align}

\paragraph{Step 5: Convert TV bounds to AUC/PR-AUC/F1 bounds.}

\emph{AUC.} The AUC can be written as
$\operatorname{AUC} = \mathbb{P}(\text{score}_{\text{noise}} >
\text{score}_{\text{clean}})$, which involves two independent samples: one
from $P(\text{score} \mid Z=1)$ (noise) and one from
$P(\text{score} \mid Z=0)$ (clean). For product measures,
\begin{align*}
&\operatorname{TV}\bigl(P(\cdot\mid Z=1) \times P(\cdot\mid Z=0),\;
        \tilde{P}(\cdot\mid Z=1) \times \tilde{P}(\cdot\mid Z=0)\bigr) \\
&\qquad\leq \operatorname{TV}(P(\cdot\mid Z=1), \tilde{P}(\cdot\mid Z=1))
         + \operatorname{TV}(P(\cdot\mid Z=0), \tilde{P}(\cdot\mid Z=0)).
\end{align*}

We transfer the marginal TV bound to the conditional distributions. For any
event $A$,
\begin{align*}
|P(A \mid Z=1) - \tilde{P}(A \mid Z=1)|
&= \frac{|P(A \cap \{Z=1\}) - \tilde{P}(A \cap \{Z=1\})|}{\eta} \\
&\leq \frac{\operatorname{TV}(P, \tilde{P})}{\eta}
\;\leq\; \frac{1}{\eta}\sqrt{\frac{\delta}{2}}.
\end{align*}
Similarly,
$\operatorname{TV}(P(\cdot\mid Z=0), \tilde{P}(\cdot\mid Z=0))
\leq \frac{1}{1-\eta}\sqrt{\frac{\delta}{2}}$.
Combining,
\begin{align*}
|\operatorname{AUC}_P(h) - \operatorname{AUC}_{\tilde{P}}(h)|
&\leq \operatorname{TV}(P(\cdot\mid Z=1), \tilde{P}(\cdot\mid Z=1))
   + \operatorname{TV}(P(\cdot\mid Z=0), \tilde{P}(\cdot\mid Z=0)) \\
&\leq \sqrt{\frac{\delta}{2}}\cdot
   \left(\frac{1}{\eta} + \frac{1}{1-\eta}\right).
\label{eq:auc-tv-transfer}
\end{align*}

\emph{PR-AUC.} \textbf{[Claim -- proof incomplete.]}
The PR-AUC involves an integral over thresholds of precision,
$\mathbb{P}(Z=1 \mid \hat{z}=1)$. The same amplification argument would apply at
each threshold, suggesting the identical bound
\eqref{eq:prauc-bound}. However, a fully rigorous treatment requires
additional steps to handle the integral over thresholds and the
$\eta$-dependent amplification of per-threshold TV bounds; the
argument presented here is therefore incomplete. We state
\eqref{eq:prauc-bound} as a claim pending a complete proof.

\emph{F1.} The F1 score is a function of the joint distribution
$P(\hat{z}, Z)$ only:
\[
\operatorname{F1} = \frac{2 \cdot \operatorname{TP}}
                         {2 \cdot \operatorname{TP}
                          + \operatorname{FP} + \operatorname{FN}},
\]
where $\operatorname{TP} = \mathbb{P}(\hat{z}=1, Z=1)$,
$\operatorname{FP} = \mathbb{P}(\hat{z}=1, Z=0)$, and
$\operatorname{FN} = \mathbb{P}(\hat{z}=0, Z=1)$. Since F1 does not involve
conditional sampling, the TV bound applies directly without amplification.

The F1 score is Lipschitz continuous in $( \operatorname{TP},
\operatorname{FP}, \operatorname{FN})$. Its partial derivatives satisfy
\[
\Bigl|\frac{\partial \operatorname{F1}}{\partial \operatorname{TP}}\Bigr|
= \frac{2(\operatorname{FP}+\operatorname{FN})}
       {(2\operatorname{TP}+\operatorname{FP}+\operatorname{FN})^2}
\leq \frac{2}{\min(\operatorname{TP}, 1)},
\]
and similarly for $\operatorname{FP}$, $\operatorname{FN}$. Let
$p_{\min} = \min(2\operatorname{TP}+\operatorname{FP}+\operatorname{FN})$.
The Lipschitz constant satisfies $C_F \leq 2/p_{\min}^2$.

In the typical operating range where precision and recall are both at least
$0.1$, direct calculation gives $C_F \leq 2$. More conservatively:
\begin{itemize}
  \item Precision, Recall $\geq 0.1$: $C_F \leq 3$,
  \item Precision, Recall $\geq 0.5$: $C_F \leq 1$.
\end{itemize}
Hence
\begin{equation}
|\operatorname{F1}_P(h) - \operatorname{F1}_{\tilde{P}}(h)|
\;\leq\; C_F \cdot \operatorname{TV}(P_{\text{pred}}, \tilde{P}_{\text{pred}})
\;\leq\; C_F \cdot \sqrt{\frac{\delta}{2}}.
\label{eq:f1-tv-transfer}
\end{equation}

\paragraph{Step 6: Combine the bounds.}
From \eqref{eq:base-equiv-au} and \eqref{eq:auc-tv-transfer},
\begin{align*}
\operatorname{AUC}_P(h_{\text{SCX}})
&\leq \operatorname{AUC}_{\tilde{P}}(h_{\text{SCX}})
   + \sqrt{\frac{\delta}{2}}\cdot\left(\frac{1}{\eta} + \frac{1}{1-\eta}\right) \\
&= \operatorname{AUC}_{\text{base}}
   + \sqrt{\frac{\delta}{2}}\cdot\left(\frac{1}{\eta} + \frac{1}{1-\eta}\right).
\end{align*}
Similarly,
\begin{align*}
\operatorname{PRAUC}_P(h_{\text{SCX}})
&\leq \operatorname{PRAUC}_{\text{base}}
   + \sqrt{\frac{\delta}{2}}\cdot\left(\frac{1}{\eta} + \frac{1}{1-\eta}\right), \\
\operatorname{F1}_P(h_{\text{SCX}})
&\leq \operatorname{F1}_{\text{base}} + C_F \cdot \sqrt{\frac{\delta}{2}}.
\end{align*}
This completes the proof.
\end{proof}

\begin{remark}[Symmetric lower bound]
\label{remark:symmetric-bound}
The theorem also admits a symmetric lower bound
\[
\operatorname{AUC}(h_{\text{SCX}}) \;\geq\;
\operatorname{AUC}_{\text{base}} -
\sqrt{\frac{\delta}{2}}\cdot\left(\frac{1}{\eta} + \frac{1}{1-\eta}\right),
\]
which follows from the same TV argument applied in the opposite direction.
When $\delta = 0$, the SCX AUC is \emph{exactly} equal to the loss baseline
AUC -- SCX can be neither better nor worse.
\end{remark}

\begin{remark}[Interpretation of the $\eta$ dependence]
\label{remark:eta-interpretation}
When $\eta$ is small (rare noise), the AUC and PR-AUC bounds become loose.
This reflects the intrinsic difficulty of detecting rare noise: SCX requires
stronger features (smaller $\delta$) to reliably surpass the loss baseline.
For $\eta = 0.5$ the factor $1/\eta + 1/(1-\eta) = 4$, the most favourable
case; for $\eta \to 0$ or $\eta \to 1$ the factor diverges.
\end{remark}

%------------------------------------------------------------------------------
\subsection{Corollaries}
\label{sec:thm2:corollaries}

\subsubsection*{Corollary 1: Complete Failure ($\delta = 0$)}
\label{cor:complete-failure}
If $\phi(X)$ is independent of $S$ (i.e.\ $\phi$ is completely uninformative),
then
\[
\operatorname{AUC}(h_{\text{SCX}}) = \operatorname{AUC}_{\text{base}},\qquad
\operatorname{F1}(h_{\text{SCX}}) = \operatorname{F1}_{\text{base}}.
\]
If, in addition, the noise is loss-uninformative (e.g.\ uniform label noise),
the loss baseline itself degenerates to random guessing:
\[
\operatorname{AUC}_{\text{base}} = 0.5,\qquad
\operatorname{F1}_{\text{base}} = \frac{2\eta}{1+\eta}.
\]

\subsubsection*{Corollary 2: Loss-Uninformative Noise}
\label{cor:loss-uninformative}
Under uniform label noise \textbf{and the assumption that experts are random
guessers} (i.e., each expert predicts labels uniformly at random independent
of the input), the loss distribution is independent of the noise
indicator:
\[
P(\ell \mid Z=1) = P(\ell \mid Z=0).
\]
The loss baseline then reduces to the random baseline, and Theorem~2 gives
\begin{align}
\operatorname{AUC}(h_{\text{SCX}}) &\leq 0.5 +
\sqrt{\frac{\delta}{2}}\cdot\left(\frac{1}{\eta} + \frac{1}{1-\eta}\right),
\label{eq:cor2-auc} \\
\operatorname{PRAUC}(h_{\text{SCX}}) &\leq \eta +
\sqrt{\frac{\delta}{2}}\cdot\left(\frac{1}{\eta} + \frac{1}{1-\eta}\right),
\label{eq:cor2-prauc} \\
\operatorname{F1}(h_{\text{SCX}}) &\leq \frac{2\eta}{1+\eta} +
C_F\sqrt{\frac{\delta}{2}}.
\label{eq:cor2-f1}
\end{align}
When $\delta = 0$, the PR-AUC is bounded by $\eta$ and the F1 by
$2\eta/(1+\eta)$.

\subsubsection*{Corollary 3: Minimum Information Threshold}
\label{cor:min-info}
Suppose we require the SCX detection F1 to exceed the loss baseline by at
least $\Delta > 0$. The minimum mutual information necessary is
\begin{equation}
\delta_{\min} \;\geq\; 2\left(\frac{\Delta}{C_F}\right)^2.
\label{eq:min-info-f1}
\end{equation}
For AUC improvements,
\begin{equation}
\delta_{\min} \;\geq\; 2\bigl(\Delta_{\text{AUC}} \cdot \eta(1-\eta)\bigr)^2.
\label{eq:min-info-auc}
\end{equation}

\begin{example}
To achieve an F1 improvement of $\Delta = 0.05$ with $C_F = 2$,
\[
\delta_{\min} \;\geq\; 2\left(\frac{0.05}{2}\right)^2
= 2 \cdot 0.000625 = 1.25 \times 10^{-3} \text{ nats}.
\]
To achieve $\Delta = 0.1$ with $C_F = 2$,
\[
\delta_{\min} \;\geq\; 2\left(\frac{0.1}{2}\right)^2
= 5 \times 10^{-3} \text{ nats}.
\]
Since $\delta$ is mutual information in nats and $\delta \leq \log K_{\mathcal{S}}$,
these thresholds are modest: $\delta = 10^{-3}$ corresponds to a likelihood
ratio of $\exp(10^{-3}) \approx 1.001$ -- very little state information is
needed for a measurable improvement.
\end{example}

\subsubsection*{Corollary 4: Effect of State Count $K_{\mathcal{S}}$}
\label{cor:state-count}
As the number of true states $K_{\mathcal{S}}$ grows, maintaining the same
$\delta = I(\phi; S)$ requires progressively stronger features. The
normalised weakness $\varepsilon_\phi = \delta / \log K_{\mathcal{S}}$
determines SCX effectiveness:
\begin{itemize}
  \item $\varepsilon_\phi \approx 1$: features carry almost no state
        information; SCX degenerates.
  \item $\varepsilon_\phi \approx 0$: features carry full state information;
        SCX can be effective.
  \item $\varepsilon_\phi \approx 0.5$: state discovery is about half
        accurate; noise detection is partially effective.
\end{itemize}
Blindly increasing the number of states $K_{\mathcal{S}}$ is therefore
counterproductive without a corresponding increase in feature information
$\delta$.

%------------------------------------------------------------------------------
\subsection{Practical Diagnostics}
\label{sec:thm2:diagnostics}

The following diagnostics operationalise Theorem~2 for practitioners.

\paragraph{Diagnostic 1: Within-state consistency convergence.}
If all estimated states have near-equal consensus scores $C(\hat{s})$,
the features may be too weak:
\[
\operatorname{Var}\bigl(\{C(\hat{s})\}_{\hat{s} \in \hat{\mathcal{S}}}\bigr)
< \tau_C.
\]

\paragraph{Diagnostic 2: Random baseline comparison.}
Compute the loss baseline AUC and PR-AUC on the training set. If they are
close to the random baseline ($\operatorname{AUC} \approx 0.5$,
$\operatorname{PRAUC} \approx \eta$), the noise detection task is intrinsically
difficult, and weak features further prevent SCX from surpassing this baseline.

\paragraph{Diagnostic 3: Supervised state matching (when available).}
If ground-truth state labels exist for a subset (e.g.\ simulated data or known
batch sources), compute the adjusted Rand index (ARI) between the $\phi$-based
clustering and the true states:
\[
\operatorname{ARI}(\hat{S}, S) < 0.1 \;\Longrightarrow\; \text{features too weak}.
\]

\paragraph{Diagnostic 4: Mutual information estimation.}
Estimate $I(\phi(X); Y)$ via a Kraskov $k$-NN estimator or MINE, and compare
with $\log K_{\mathcal{S}}$:
\[
\frac{\hat{I}(\phi(X); Y)}{\log K_{\mathcal{S}}} < 0.2
\;\Longrightarrow\; \text{features may be too weak}.
\]

\paragraph{Threshold guidelines.}
Based on Theorem~2 and Corollary~3, we recommend the operational thresholds
in Table~\ref{tab:diagnostic-thresholds}.

\begin{table}[h]
\centering
\caption{Practical guidelines for feature strength assessment.}
\label{tab:diagnostic-thresholds}
\begin{tabular}{lc}
\toprule
Condition & Recommendation \\
\midrule
$\varepsilon_\phi < 0.2$            & SCX likely effective \\
$0.2 \leq \varepsilon_\phi \leq 0.5$ & SCX partially effective; check other metrics \\
$\varepsilon_\phi > 0.5$            & SCX not recommended; strengthen features first \\
$\operatorname{AUC}_{\text{base}} < 0.55$ and $\varepsilon_\phi > 0.5$
                                    & SCX noise detection will fail \\
$\operatorname{AUC}_{\text{base}} > 0.7$
                                    & Loss baseline already strong; SCX may help marginally \\
$\eta < 0.1$ and $\varepsilon_\phi > 0.3$
                                    & Rare noise amplifies AUC bound; need stronger features \\
$0.3 \leq \eta \leq 0.7$           & Moderate noise rate; AUC/PR-AUC bounds most favourable \\
\bottomrule
\end{tabular}
\end{table}

Here $\varepsilon_\phi = I(\phi; S) / \log K_{\mathcal{S}}$ and
$\eta = \mathbb{P}(Z=1)$ is the marginal noise rate.

%------------------------------------------------------------------------------
% Local bibliography entries
\begin{thebibliography}{10}
\bibitem{fano1961}
R.~M.~Fano, \emph{Transmission of Information: A Statistical Theory of
Communications}. MIT Press, 1961.

\bibitem{cover2006}
T.~M.~Cover and J.~A.~Thomas, \emph{Elements of Information Theory}, 2nd~ed.
Wiley, 2006.

\bibitem{pinsker1964}
M.~S.~Pinsker, \emph{Information and Information Stability of Random Variables
and Processes}. Holden-Day, 1964.
\end{thebibliography}

\endinput
